\section{Proofs}
\label{app:proofs}

\subsection{Proof of Proposition~\ref{prop:rank_inflation}}

Let
\[
  S_{\mathbf{i}}
  =
  \sum_{r_1,\ldots,r_N}
  X_{r_1,\ldots,r_N}
  \prod_{n=1}^{N}A^{(n)}_{i_n,r_n}.
\]
Expanding $S_{\mathbf{i}}^p$ gives products of $p$ Tucker index tuples. For
each mode $n$, the resulting mode-dependent terms are all degree-$p$ monomials
in the row $A^{(n)}_{i_n,:}$. These monomials are indexed by multi-indices
$\alpha\in\mathbb{N}^{R_n}$ with $|\alpha|=p$, of which there are
$\binom{R_n+p-1}{p}$. Collecting identical monomials yields a Tucker
representation of $\cS^{\od p}$ whose mode-$n$ factor is the degree-$p$
Veronese lift of $\bA^{(n)}$. This proves the stated rank bound.

\subsection{Proof of Theorem~\ref{thm:tucker_separation}}

We first record a standard fact about Veronese evaluations. Let
\(D=\binom{R+p-1}{p}\). There exists a matrix \(\bA\in\R^{D\times R}\) whose
degree-\(p\) Veronese feature matrix has rank \(D\). Indeed, the homogeneous
degree-\(p\) monomials in \(R\) variables are linearly independent
polynomials, so one can choose \(D\) evaluation points for which the square
evaluation matrix is nonsingular. Consequently, for \(I\ge D\), a generic
\(\bA\in\R^{I\times R}\) has a degree-\(p\) Veronese lift with full column
rank \(D\).

Choose generic factor matrices
\(\bA^{(n)}=[a_1^{(n)},\ldots,a_R^{(n)}]\in\R^{I_n\times R}\) and define the
CP-diagonal Tucker signal
\[
  \cS =
  \sum_{r=1}^{R} a_r^{(1)}\otimes\cdots\otimes a_r^{(N)} .
\]
For generic factors, each mode unfolding of \(\cS\) has rank \(R\), so the
signal has exact Tucker rank \((R,\ldots,R)\).

Now set the degree-\(p\) response coefficient to one and all other coefficients to
zero, so \(\cY=\cS^{\od p}\). Expanding entrywise gives
\[
  Y_{\mathbf{i}}
  =
  \left(
  \sum_{r=1}^{R}\prod_{n=1}^{N} A^{(n)}_{i_n,r}
  \right)^p
  =
  \sum_{|\alpha|=p}
  c_\alpha
  \prod_{n=1}^{N}
  \prod_{r=1}^{R}
  \left(A^{(n)}_{i_n,r}\right)^{\alpha_r},
\]
where \(c_\alpha=p!/(\alpha_1!\cdots\alpha_R!)>0\). Let
\(\bB^{(n)}=\bA^{(n)\langle p\rangle}\in\R^{I_n\times D}\) be the Veronese
lift. By the preceding full-rank fact, each \(\bB^{(n)}\) has column rank
\(D\) for generic \(\bA^{(n)}\). The mode-\(n\) unfolding of \(\cY\) factors
as \(\bB^{(n)}\) times a full-column-rank right factor built from the
remaining lifted factors, with a nonsingular diagonal scaling by the positive
coefficients \(c_\alpha\). Therefore every mode unfolding has rank \(D\). The
exact standard Tucker rank of \(\cY\) is \((D,\ldots,D)\), while \(\cY\) is
represented by the degree-\(p\) linked model using the rank-\((R,\ldots,R)\)
Tucker signal.

\subsection{Proof of Corollary~\ref{cor:tied_lift_parameters}}

By Proposition~\ref{prop:rank_inflation}, the degree-$p$ channel can be written
as a Tucker tensor with ranks at most
$D_n(p)=\binom{R_n+p-1}{p}$. A standalone Tucker parameterization has
$\prod_n D_n(p)$ core entries and $\sum_n I_nD_n(p)$ factor entries. In
contrast, the power-dictionary linked model reuses the core and factors of
$\cS$ and adds only one response coefficient per channel. Summing degrees up to
$Q$ adds $Q+1$ response coefficients, proving the stated comparison.

\subsection{Proof of Theorem~\ref{thm:poly_generation}}

If $g(s)=\sum_{p=0}^{P}\gamma_p s^p$ and $\cY=g(\cS^\star)$, choose the
Tucker component of the power-dictionary linked model to be the Tucker representation of
$\cS^\star$ and set $\beta_p=\gamma_p$ for all $p=0,\ldots,P$. Then
$f_{\bbeta}(\cS^\star)=g(\cS^\star)=\cY$, giving exact representation with the
same Tucker rank.

\subsection{Proof of Proposition~\ref{prop:response_approximation}}

By the Weierstrass approximation theorem, for every $\epsilon>0$ there exists
a polynomial $q(s)=\sum_{p=0}^{P}\beta_ps^p$ such that
$\sup_{s\in[a,b]}|f(s)-q(s)|\le\epsilon$. Every entry of $\cS^\star$ lies in
$[a,b]$, so each entrywise approximation error is bounded by $\epsilon$.
Therefore
\[
  \norm{f(\cS^\star)-q(\cS^\star)}_F^2
  =
  \sum_{\mathbf{i}}
  |f(S^\star_{\mathbf{i}})-q(S^\star_{\mathbf{i}})|^2
  \le
  M\epsilon^2,
\]
which gives the stated bound after taking square roots.

\subsection{Proof of Proposition~\ref{prop:monotone_descent}}

For a fixed Tucker signal, Response Refresh solves the $\bbeta$-subproblem
exactly, hence it cannot increase $\mathcal{F}_Q$. The subsequent Tucker
update is assumed to satisfy the sufficient decrease condition
\eqref{eq:sufficient_decrease}, and therefore also cannot increase
$\mathcal{F}_Q$. Combining the two stages gives monotone decrease over the
outer iteration. Since $\mathcal{F}_Q$ is bounded below by zero when the
regularization weights are nonnegative, summing
\eqref{eq:sufficient_decrease} over $k$ gives
\[
  c\sum_{k=0}^{K}\norm{\theta^{k+1}-\theta^k}_2^2
  \le
  \mathcal{F}_Q(\theta^0)-\mathcal{F}_Q(\theta^{K+1})
  \le
  \mathcal{F}_Q(\theta^0).
\]
Letting $K\to\infty$ proves square summability of the steps.

\subsection{Proof of Theorem~\ref{thm:stationary_convergence}}

For fixed $Q$, the objective is a polynomial function of the Tucker parameters
and response coefficients when the monomial basis is used. For any other fixed
polynomial basis spanning the same space, it remains polynomial after a linear
change of coefficients. The linear-spline dictionary is piecewise polynomial
and semi-algebraic; standard RBF dictionaries are analytic and belong to the
standard tame classes covered by KL convergence theory. With
standard box or normalization constraints used to fix the Tucker gauge, the
corresponding constrained objective satisfies the Kurdyka--Lojasiewicz
property. Assumption~\ref{ass:optimization} gives sufficient decrease and a
relative error bound. The standard KL convergence theorem for nonconvex descent
sequences then implies finite length of the iterates and convergence to a
critical point of $\mathcal{F}_Q$.

\subsection{Proof of Theorem~\ref{thm:masked_generalization}}

Fix \(0<\varepsilon\le B_\Theta\), and let
\(\mathcal{N}_\varepsilon\) be an \(\varepsilon\)-net of \(\Theta\) in
Euclidean norm. Since \(\Theta\) lies in a \(d\)-dimensional ball of radius
\(B_\Theta\), it admits a net with
\[
  |\mathcal{N}_\varepsilon|
  \le
  \left(\frac{3B_\Theta}{\varepsilon}\right)^d .
\]
For any fixed \(\theta_0\in\mathcal{N}_\varepsilon\), the random variables
\(\ell_{\mathbf{i}_t}(\theta_0)\) are independent and bounded in
\([0,B_\ell]\). Hoeffding's inequality gives
\[
  \Pr\left(
  \left|
  \mathcal{L}(\theta_0)-\widehat{\mathcal{L}}_m(\theta_0)
  \right|
  \ge u
  \right)
  \le
  2\exp\left(-\frac{2mu^2}{B_\ell^2}\right).
\]
A union bound over \(\mathcal{N}_\varepsilon\) shows that, with probability at
least \(1-\delta\),
\[
  \max_{\theta_0\in\mathcal{N}_\varepsilon}
  \left|
  \mathcal{L}(\theta_0)-\widehat{\mathcal{L}}_m(\theta_0)
  \right|
  \le
  B_\ell
  \sqrt{
  \frac{
  d\log(3B_\Theta/\varepsilon)+\log(2/\delta)}
  {2m}} .
\]
For an arbitrary \(\theta\in\Theta\), choose
\(\theta_0\in\mathcal{N}_\varepsilon\) with
\(\norm{\theta-\theta_0}_2\le\varepsilon\). The uniform Lipschitz assumption
implies
\[
  |\mathcal{L}(\theta)-\mathcal{L}(\theta_0)|\le G\varepsilon,
  \qquad
  |\widehat{\mathcal{L}}_m(\theta)-\widehat{\mathcal{L}}_m(\theta_0)|
  \le G\varepsilon .
\]
Combining these two approximation errors with the net bound and taking the
supremum over \(\Theta\) gives the stated inequality.
